\begin{table}[htp]
\centering

    % Esta primera línea define las columnas de la tabla. Los posibles tipos de columna son:
    % c: texto centrado
    % l: texto alineado a la izquierda
    % r: texto centrado a la derecha
    % p: columna de ancho fijo
    % Las columnas tienen ancho dinámico según la anchura máxima de los elementos que contengan.

    % Las columnas l/r/c no parten el texto en filas diferentes si éste es demasiado largo. Para ello, puede utilizar el tipo de columna de ancho fijo "p".
    
    % Las barras verticales | se usan para definir los bordes verticales de la tabla. Pruebe a eliminar algunas y observe qué ocurre.
    \begin{center}
    \begin{tabular}{ | l | l | }
        
        % A continuación van las filas de la tabla. En cada fila, las columnas se separan con el carácter &
        % Para terminar una fila se usa \\
        % Para incluir un borde horizontal entre filas se usa \hline
        
        % Cabecera con textos en negrita:
        \hline
        \textbf{Media: }  & 0.206815 \\
        \hline
        \textbf{Mediana} & 0.192512 \\
        \hline
        \textbf{Moda} & 0.163906 \\
        \hline
        \textbf{IQR} & 0.047104 \\
        \hline
        \textbf{Varianza} & 0.002715 \\
        \hline
        \textbf{Desviación estándar} & 0.052101 \\
        \hline
        \textbf{Covarianza (GPU)} & 0.000971 \\
        \hline
        \textbf{Correlación (GPU)} & 0.419512 \\
        \hline
        \textbf{Q1} & 0.190464 \\
        \hline
        \textbf{Q2} & 0.192512 \\
        \hline
        \textbf{Q3} & 0.237568 \\
        \hline
        
    \end{tabular} 
    \end{center}
    
    \caption{Resultados de la renderización en la prueba del polígono: \textbf{GPU}}
    
    \label{table:results-animation-render-gpu}
\end{table}